% ======================================================================
% Přednáška 11 (12. 12. 2022)
% Shrnutí neutrinových experimentů, CP narušení, Dirac x Majorana,
% supersymetrie, urychlovače částic
% ======================================================================
\section{Otevřené otázky neutrin, supersymetrie a urychlovače částic}

\strana{P11-1}
\subsection{Shrnutí neutrinových experimentů}

Na minulých přednáškách jsme probrali oscilace neutrin. Experimenty, které je změřily, lze
rozdělit podle zdroje neutrin -- každý zdroj dává jinou vůni, jinou energii \(E\) a jinou
vzdálenost \(L\) mezi zdrojem a detektorem, a je proto citlivý na jiný parametr oscilací
(fáze oscilace je úměrná \(\Delta m^2 L/E\)):
\begin{itemize}[nosep]
  \item \kw{solární} (\(\nu_e\)) -- z termojaderných reakcí ve Slunci (\(E\sim\mathrm{MeV}\),
        \(L\approx 1{,}5\cdot10^{11}\,\mathrm{m}\)); měří \(\theta_{12}\), \(\Delta m^2_{21}\)
        (Homestake, SNO, Super-Kamiokande\dots);
  \item \kw{reaktorová} (\(\bar\nu_e\)) -- z \(\beta^-\) rozpadů štěpných produktů v jaderném
        reaktoru (\(E\sim\) jednotky MeV); měří zánik \(\bar\nu_e\), např. KamLAND, Daya Bay
        (\(\theta_{13}\));
  \item \kw{atmosférická} -- z rozpadů pionů a mionů vzniklých interakcí kosmického záření v
        atmosféře (\(\nu_\mu,\bar\nu_\mu,\nu_e,\bar\nu_e\)); \(L\) od \(\sim 15\,\mathrm{km}\)
        (shora) do \(\sim 13\,000\,\mathrm{km}\) (zespodu, přes celou Zemi); Super-Kamiokande
        zde v roce 1998 oscilace objevil (\(\theta_{23}\), \(|\Delta m^2_{32}|\));
  \item \kw{urychlovačová} -- uměle vyrobené svazky \(\nu_\mu\) nebo \(\bar\nu_\mu\) se známou energií
        a vzdáleností (T2K, NOvA, OPERA\dots).
\end{itemize}

\subsection{CP narušení u neutrin}

Proč nás zajímá CP narušení? Vesmír obsahuje téměř výhradně hmotu a ne antihmotu; jednou z
podmínek, aby taková asymetrie mohla vzniknout, je narušení symetrie \(CP\) (Sacharovovy
podmínky). V kvarkovém sektoru CP narušení známe (kaony, \(B\)-mezony), ale je příliš malé.
Neutrinová směšovací (PMNS) matice obsahuje obecně fázi \(\delta\), která CP narušuje -- a
oscilace neutrin jsou místem, kde ji lze změřit.

Operace \(CP\) převádí neutrino na antineutrino (\(C\)) a obrací helicitu (\(P\)), takže z
levotočivého \(\nu\) udělá pravotočivé \(\bar\nu\) -- tedy opět fyzikální částici. Porovnáváme
proto oscilace neutrin a antineutrin:
\begin{itemize}
  \item \kw{zachování CP}:
        \[
          P_{\nu_\alpha\to\nu_\beta} = P_{\bar\nu_\alpha\to\bar\nu_\beta},
        \]
        kde \(P_{\nu_\alpha\to\nu_\beta}\) je pravděpodobnost, že neutrino vzniklé s vůní
        \(\alpha\) bude detekováno s vůní \(\beta\) (\(\alpha,\beta\in\{e,\mu,\tau\}\));
  \item \kw{narušení CP} (\(\not{\!CP}\)):
        \[
          \vzorec{P_{\nu_\alpha\to\nu_\beta} - P_{\nu_\beta\to\nu_\alpha}
          = P_{\nu_\alpha\to\nu_\beta} - P_{\bar\nu_\alpha\to\bar\nu_\beta}
          \;\propto\; 2\sin\delta}
        \]
\end{itemize}
První rozdíl (záměna počátečního a koncového stavu) měří narušení \emph{časové}
symetrie \(T\). Rovnost s druhým rozdílem plyne z věty \(CPT\): \(CPT\) symetrie zaručuje
\(P_{\nu_\beta\to\nu_\alpha} = P_{\bar\nu_\alpha\to\bar\nu_\beta}\), takže narušení \(T\) je
ekvivalentní narušení \(CP\). Velikost asymetrie je úměrná \(\sin\delta\); je-li \(\delta=0\)
nebo \(\pi\), CP narušeno není.

\begin{vysvetleni}[Proč je potřeba tři vůně]
  Při oscilacích jen dvou neutrin lze matici směšování zvolit reálnou a CP fáze neexistuje.
  Pro tři generace je asymetrie (ve vakuu)
  \(P_{\nu_\mu\to\nu_e}-P_{\bar\nu_\mu\to\bar\nu_e} \propto J\,
  \sin\Delta_{21}\sin\Delta_{32}\sin\Delta_{31}\), kde \(\Delta_{ij}=\Delta m^2_{ij}L/(4\hbar cE)\)
  a Jarlskogův invariant \(J=c_{12}s_{12}c_{23}s_{23}c_{13}^2s_{13}\sin\delta\). Asymetrie je
  tedy nenulová jen tehdy, když jsou \emph{všechny} tři směšovací úhly nenulové (proto bylo
  měření \(\theta_{13}\neq0\) v roce 2012 tak důležité). Pozor: při průletu Zemí vzniká i
  ,,falešná`` asymetrie od interakce s elektrony v látce (MSW jev), kterou je třeba odečíst.
\end{vysvetleni}

\strana{P11-2}
Jak se to měří na urychlovači? Svazek neutrin vzniká z rozpadů pionů,
\[
  \pi^+\to\mu^+\nu_\mu,\qquad \pi^-\to\mu^-\bar\nu_\mu .
\]
Magnetickým ,,rohem`` (horn) se fokusují buď kladné, nebo záporné piony -- přepnutím
polarity tak dostaneme buď svazek \(\nu_\mu\), nebo \(\bar\nu_\mu\), a ve vzdáleném detektoru
porovnáváme \(\nu_\mu\to\nu_e\) a \(\bar\nu_\mu\to\bar\nu_e\) (T2K, NOvA, budoucí DUNE).

Pro tauonová neutrina připomeňme zdroj z minulé přednášky:
\[
  D_s^+ = (c\bar s),\qquad D_s^+\to\tau^+\nu_\tau \quad (4\,\%)
\]
\pozn{v poznámkách ,,(4\,\%)``, příp. ,,7\,\%``; PDG uvádí \(\approx5{,}4\,\%\)}

\subsection{Dirac \texorpdfstring{$\times$}{x} Majorana}

Neutrino je jediný fermion bez elektrického náboje. Otevírá to otázku, zda je neutrino
odlišné od antineutrina (\kw{Diracova částice}, jako elektron a pozitron), nebo zda je
neutrino \emph{svou vlastní antičásticí} (\kw{Majoranova částice}); v tom případě by se
,,neutrino`` a ,,antineutrino`` lišily jen helicitou. Rozhodnout to lze pomocí
\kw{dvojitého beta rozpadu}.

Obyčejný \(\beta^-\) rozpad zvýší náboj jádra o jedna:
\begin{align*}
  {}^{A}_{Z}X &\to {}^{A}_{Z+1}Y + e^- + \bar\nu_e ,\\
  {}^{A}_{Z+1}Y &\to {}^{A}_{Z+2}G + e^- + \bar\nu_e .
\end{align*}
Normálně by tedy jádro přešlo na \(Z+2\) dvěma po sobě jdoucími rozpady. U jader se sudým
\(A\) to ale nemusí jít. Z Weizsäckerovy formule plyne, že pro dané \(A\) leží hmoty jader
\(M(A,Z)\) jako funkce \(Z\) na parabole; pro sudé \(A\) jsou to dvě paraboly, protože párovací
člen \(\delta\) zvýhodňuje jádra \kw{sudo-sudá} (SS, sudé \(Z\) i \(N\)) a znevýhodňuje
\kw{licho-lichá} (LL):

\begin{center}
\begin{tikzpicture}[>=Stealth,xscale=0.9,yscale=0.85]
  \draw[thick,->] (0,0) -- (8.6,0) node[below] {$Z$};
  \draw[thick,->] (0,0) -- (0,5.2) node[left] {$M(A,Z)$};
  % SS parabola
  \draw[thick,domain=0.3:8.0,samples=60] plot (\x,{0.2+0.16*(\x-4)^2});
  \node[right] at (8.0,2.9) {SS};
  % LL parabola (carkovane)
  \draw[thick,dashed,domain=0.8:8.0,samples=60] plot (\x,{1.6+0.16*(\x-4)^2});
  \node[right] at (8.0,4.3) {LL};
  % body: SS v Z=1,3,5,7 ; LL v Z=2,4,6
  \foreach \z in {1,3,5,7} \fill ({\z},{0.2+0.16*(\z-4)^2}) circle (2.5pt);
  \foreach \z in {2,4,6} \fill ({\z},{1.6+0.16*(\z-4)^2}) circle (2.5pt);
  % prechody
  \draw[ForestGreen,thick,->] (2,2.24) -- (1.08,1.68) node[midway,above left] {$\beta^+$};
  \draw[ForestGreen,thick,->] (2,2.24) -- (2.95,0.42) node[pos=0.4,right] {$\beta^-$};
  \draw[ForestGreen,thick,dashed,->] (3,0.36) -- (3.95,1.55);
  \draw[ForestGreen,thick] (3.3,1.3) -- (3.7,0.7);
  \draw[ForestGreen,thick] (3.3,0.7) -- (3.7,1.3);
  \draw[ForestGreen,thick] (3,0.36) circle (0.25);
  \node[ForestGreen] at (2.5,-0.05) {\large ?};
  \draw[ForestGreen,thick,->] (3.15,0.15) to[bend right=10] (4.9,0.1);
  \node[ForestGreen,below] at (4.4,-0.05) {$2\beta^-$};
\end{tikzpicture}
\end{center}

Kroužkem označené sudo-sudé jádro \((Z)\) je lehčí než sousední licho-liché jádro \((Z+1)\),
takže obyčejný \(\beta^-\) rozpad je energeticky zakázán (přeškrtnutá šipka). Sudo-sudé jádro
\((Z+2)\) je ale ještě lehčí. Jediná cesta je přeskočit mezistav a změnit náboj rovnou o 2 --
\kw{dvojitý \(\beta\) rozpad}, proces druhého řádu slabé interakce s extrémně dlouhými
poločasy (\(\sim10^{19}\)--\(10^{24}\) let; např. \({}^{76}\mathrm{Ge}\), \({}^{130}\mathrm{Te}\),
\({}^{136}\mathrm{Xe}\)).

\strana{P11-3}
Možné jsou dva typy:
\begin{align*}
  \text{Dirac:}\qquad & {}^{A}_{Z}X \to {}^{A}_{Z+2}Y + e^- + e^- + \bar\nu_e + \bar\nu_e
     && (2\nu\beta\beta),\\
  \text{Majorana:}\qquad & {}^{A}_{Z}X \to {}^{A}_{Z+2}Y + e^- + e^- && (0\nu\beta\beta).
\end{align*}
Dvouneutrinový rozpad \(2\nu\beta\beta\) je prostě ,,dva \(\beta\) rozpady najednou`` a je povolen
v každém případě (pozorován u řady jader). Bezneutrinový rozpad \(0\nu\beta\beta\) je možný jen
tehdy, je-li neutrino Majoranova částice: antineutrino vyzářené v jednom vrcholu se v druhém
vrcholu pohltí jako neutrino (\(\bar\nu_e\equiv\nu_e\)). Navíc musí neutrino mít nenulovou
hmotu, aby se mohla ,,překlopit`` jeho helicita (amplituda je úměrná efektivní Majoranově
hmotě \(m_{\beta\beta}\)). Proces narušuje leptonové číslo o 2.

\begin{vysvetleni}[Jak se \(0\nu\beta\beta\) pozná]
  Experimentálně se měří \emph{součet energií obou elektronů}. Při \(2\nu\beta\beta\) odnášejí
  část energie neutrina, takže spektrum je spojité od nuly do \(Q_{\beta\beta}\). Při
  \(0\nu\beta\beta\) odnesou elektrony celou energii rozpadu (zpětný ráz jádra je zanedbatelný),
  a ve spektru by se tak objevil úzký pík přesně na \(Q_{\beta\beta}\). Zatím pozorován nebyl,
  limity na poločas jsou \(T_{1/2}\gtrsim10^{26}\) let (GERDA, KamLAND-Zen), viz Davídkovy slajdy
  ,,Otevřené otázky částicové fyziky``.
\end{vysvetleni}

\subsection{Supersymetrie?}

Další otevřenou otázkou (,,co ještě nevíme``) je, zda existuje \kw{supersymetrie} (SUSY) --
symetrie mezi fermiony a bosony. Motivace: stabilizuje hmotu Higgsova bosonu vůči velkým
kvantovým korekcím, zlepšuje sjednocení vazbových konstant při vysokých energiích a nejlehčí
supersymetrická částice je kandidátem na temnou hmotu.
\begin{itemize}
  \item Každý fermion má ,,kamaráda`` -- \kw{skalární částici} (\(s=0\)), a to \(2\times\)
        (jednu pro každou chiralitu):
        \[
          \mu \to \tilde\mu_L,\ \tilde\mu_R ,
        \]
        a stejně pro (ostatní) leptony; obdobně pro kvarky (skvarky \(\tilde q_L,\tilde q_R\)).
  \item Obdobně pro \(g,\gamma,W,Z\): bosony dostanou partnery se spinem \(\tfrac12\)
        (gluino \(\tilde g\), \dots).
  \item Higgsový sektor se rozšíří na \(5\times\) Higgs: \(H^\pm,\ h^0,\ H^0,\ A^0\)
        (\(H^+\) a \(H^-\) jsou dva stavy, celkem tedy pět).
\end{itemize}
Partneři elektroslabých bosonů a Higgsů se navzájem mísí (mají stejná kvantová čísla), takže
fyzikální stavy jsou jejich kombinace:
\begin{align*}
  \gamma,\,Z\ (\text{a neutrální Higgsy})&:\quad \tilde\chi^0_1,\ \tilde\chi^0_2,\ \dots
     && s=\tfrac12 \quad\text{(čtyři \kw{neutralina})},\\
  W^\pm,\,H^\pm&:\quad \tilde\chi^\pm_1,\ \tilde\chi^\pm_2,\ \dots && s=\tfrac12
     \quad\text{(dvě \kw{chargina})}.
\end{align*}

\strana{P11-4}
\textbf{Jak se hledají:} Supersymetrické částice by na LHC vznikaly v párech (zachování
\(R\)-parity), např. pár skvarků, který se rozpadá
\[
  \tilde q\,\tilde q \to q\,q\,\tilde\chi^0_1\,\tilde\chi^0_2 .
\]
Neutralina jsou neutrální a slabě interagující; nejlehčí z nich (\(\tilde\chi^0_1\), LSP) je
stabilní a detektorem projde jako neutrino (těžší \(\tilde\chi^0_2\) se dál rozpadá až na
\(\tilde\chi^0_1\)). Pozorujeme proto
\[
  \text{2 jety} + \text{velkou chybějící } p_T .
\]
Tato signatura je stejná jako u některých procesů SM (např. \(Z\to\nu\bar\nu\) + jety), hledá
se tedy přebytek případů s velmi velkou chybějící \(p_T\). Zatím nic nalezeno nebylo a hmoty
skvarků a gluin pod \(\sim 1\)--\(2\,\mathrm{TeV}\) jsou vyloučeny.

\subsection{Urychlovače částic}

Poslední velké téma: jak částice na potřebné energie urychlit. Urychlovat lze jen nabité
částice, a to jedině elektrickým polem (magnetické pole mění jen směr). Dělíme je na
\[
  \text{\kw{lineární}} \quad\times\quad \text{\kw{kruhové}} .
\]
Lineární urychlovače dále dělíme na
\begin{itemize}[nosep]
  \item \kw{elektrostatické} -- např. Van de Graaffův generátor (VdG);
  \item \kw{rezonanční} -- urychlují střídavým vysokofrekvenčním polem.
\end{itemize}

\subsubsection{Elektrostatické urychlovače}

\begin{center}
\begin{tikzpicture}[>=Stealth,scale=0.8]
  % kopule
  \draw[thick] (0.2,5.0) arc[start angle=200,end angle=-20,x radius=2.0,y radius=1.6];
  \draw[thick] (0.2,5.0) -- (0.5,4.4);
  \foreach \p in {(0.0,5.6),(0.3,6.6),(1.2,7.2),(2.0,7.35),(2.8,7.2),(3.6,6.7),(4.0,5.7),(3.8,4.6)}
     \node at \p {\small $+$};
  % pas
  \draw[thick] (2.0,3.5) ellipse (0.55 and 2.5);
  \draw[thick] (2.0,5.5) circle (0.3);
  \draw[thick] (2.0,1.5) circle (0.3);
  \draw[->] (2.15,4.5) arc[start angle=120,end angle=-60,radius=0.2];
  % hrebinek nahore
  \draw[thick,->] (2.0,5.8) -- (0.4,5.7);
  % trubice svazku
  \draw[thick] (3.6,4.65) -- (3.6,0.5);
  \draw[thick] (3.85,4.65) -- (3.85,0.5);
  \draw[ForestGreen,thick,->] (3.72,4.0) -- (3.72,2.4);
  \node[right] at (4.0,2.8) {$X^+$};
  % uzemneni
  \draw[thick] (3.2,0.5) -- (4.2,0.5);
  \draw[thick] (3.4,0.3) -- (4.0,0.3);
  \draw[thick] (3.6,0.1) -- (3.8,0.1);
  % HV
  \draw[thick] (-0.6,0.2) rectangle (0.3,0.9); \node at (-0.15,0.55) {HV};
  \draw[thick,->] (-0.15,0.9) -- (-0.15,1.4) -- (1.35,1.5);
  \node at (0.5,1.75) {\small $+\!+\!+$};
  % vstup X-
  \draw[thick,->] (3.7,10.0) -- (3.7,8.0);
  \node[right] at (3.8,9.0) {$X^-$};
  \node[ForestGreen,right] at (5.2,7.0) {elektrostatický};
  \node[right] at (5.2,3.2) {speciální konfigurace:};
  \node[right] at (5.2,2.5) {tandem};
\end{tikzpicture}
\end{center}
\pozn{v náčrtu z poznámek vstupuje svazek \(X^-\) shora do kopule a \(X^+\) vylétá trubicí
dolů k uzemnění}

\kw{Van de Graaffův generátor}: zdroj vysokého napětí (HV) ,,nastříká`` kladný náboj na
nekonečný izolační pás. Pás (poháněný kladkami) náboj mechanicky vynáší nahoru do duté kovové
kopule, kde ho hrot (hřebínek) sejme; náboj přejde na vnější povrch kopule (uvnitř vodiče je
pole nulové, takže lze stále přidávat další náboj). Kopule se tak nabíjí na napětí \(U\) až
\(\sim\)\,MV (max.\ asi \(25\,\mathrm{MV}\), limituje to elektrický průraz; proto bývá celé
zařízení v tlakové nádobě s izolačním plynem \(\mathrm{SF_6}\)). Iont s nábojem \(Qe\) urychlený
mezi kopulí a uzemněním získá energii \(E = QeU\). Výhodou je spojitý svazek s velmi přesně
definovanou energií.

\kw{Tandem} (speciální konfigurace): Záporné ionty \(X^-\) se vstřelí z uzemněné strany a
urychlují se \emph{k} kladné kopuli (zisk \(eU\)). V kopuli projdou tenkou fólií nebo plynem
(,,stripper``), který jim strhne elektrony, takže se z nich stanou kladné ionty \(X^{Q+}\).
Ty jsou nyní od kopule odpuzovány a urychlují se podruhé, zpět k zemi (zisk \(QeU\)). Celkem
\[
  E = (1+Q)\,eU ,
\]
tedy stejné napětí je využito dvakrát (a vícekrát pro vícenásobně ionizované ionty) a zdroj
iontů i terč jsou na zemním potenciálu.

\strana{P11-5}
\subsubsection{Rezonanční (lineární) urychlovače}

Elektrostatické napětí nelze zvyšovat donekonečna (průraz). Řešení: použít \emph{malé}
střídavé napětí mnohokrát za sebou. Nejjednodušší je \kw{Wideröeho urychlovač} s driftovými
trubicemi:

\begin{center}
\begin{tikzpicture}[>=Stealth]
  % zdroj
  \draw[thick] (0.3,1.4) -- (0,1.4) -- (0,0.6) -- (0.3,0.6);
  \draw[thick] (0.1,0.7) rectangle (0.25,1.3);
  \draw[thick] (0,1.0) -- (7.6,1.0);
  % trubice
  \draw[ForestGreen,thick] (1.2,0.75) rectangle (2.0,1.25);
  \draw[ForestGreen,thick] (2.8,0.75) rectangle (4.0,1.25);
  \draw[ForestGreen,thick] (4.8,0.7) rectangle (6.6,1.3);
  \node at (1.7,0.85) {\small $e^-$};
  \node at (0.5,1.75) {$e^-$};
  \node at (1.5,1.9) {$+$}; \node at (1.5,1.5) {$-$};
  \node at (3.4,1.6) {$+$};
  % spoje
  \draw[thick] (0.05,0.6) -- (0.05,-0.3);
  \draw[thick] (3.4,0.75) -- (3.4,-0.3);
  \fill (3.4,-0.3) circle (1.5pt);
  \draw[thick] (0.05,-0.3) -- (6.8,-0.3) -- (6.8,-0.6);
  \fill (0.05,-0.3) circle (1.5pt);
  \draw[thick] (0.05,-0.3) -- (0.05,-0.6);
  \draw (-0.15,-0.6) -- (0.25,-0.6); \draw (-0.08,-0.7) -- (0.18,-0.7); \draw (0,-0.8) -- (0.1,-0.8);
  \draw[thick] (1.6,0.75) -- (1.6,-1.0) -- (5.7,-1.0);
  \draw[thick] (5.7,0.7) -- (5.7,-1.0);
  \fill (5.7,-1.0) circle (1.5pt);
  \draw[thick] (5.7,-1.0) -- (6.0,-1.0);
  % RF generator
  \draw[thick] (6.0,-1.4) rectangle (7.2,-0.6);
  \draw (6.6,-1.0) circle (0.25); \node at (6.6,-1.0) {\small $\sim$};
  \node[right] at (7.3,-1.0) {RF};
\end{tikzpicture}
\end{center}

Trubice jsou střídavě připojeny k pólům vysokofrekvenčního (RF) generátoru o frekvenci
\(f\). Uvnitř trubice je pole nulové (Faradayova klec) -- částice tam jen volně letí
(,,driftuje``). Urychlí se vždy v mezeře mezi trubicemi. Trik je v tom, že doba průletu
trubicí je přesně půl periody: když částice vyletí z trubice, napětí se mezitím otočilo a
v další mezeře ji opět urychlí. Pokud částice projde každou mezerou ve fázi \(\varphi_0\)
napětí \(U(t)=U_{\max}\sin\varphi\), získá v každé mezeře energii \(QeU_{\max}\sin\varphi_0\) a
po \(i\) mezerách má (nerelativisticky) kinetickou energii
\[
  E_i = i\,Q e\,U_{\max}\sin\varphi_0 = \tfrac12 m v_i^2 .
\]
Délka \(i\)-té trubice musí odpovídat dráze, kterou částice urazí za půl periody
\(T/2 = 1/(2f)\):
\[
  \vzorec{l_i = \frac{v_i}{2f}
  = \frac{1}{2f}\sqrt{\frac{2iQeU_{\max}\sin\varphi_0}{m}}
  = \frac{1}{f}\sqrt{\frac{iQeU_{\max}\sin\varphi_0}{2m}}}
\]
Délky trubic tedy rostou jako \(\sqrt i\). Pro relativistické částice (\(v\to c\)) se délka
ustálí na \(l=c/(2f)\) -- proto se tento typ hodí hlavně pro pomalé těžké ionty; pro
elektrony, které jsou relativistické už od \(\sim\)\,MeV, se používají vlnovody a rezonátory
s vysokou frekvencí.

\strana{P11-6}
\textbf{Fázová stabilita.} Částice ve svazku nemají přesně stejné energie, a tedy ani
přesně stejný čas příletu do mezery. Je třeba, aby se odchylky samy opravovaly:

\begin{center}
\begin{tikzpicture}[>=Stealth,xscale=1.1]
  \draw[thick,->] (-0.3,0) -- (6.6,0) node[below] {$t$};
  \draw[thick,->] (0,-1.6) -- (0,1.9) node[left] {$U$};
  \draw[thick,domain=0:6.2,samples=80] plot (\x,{1.4*sin(\x/6.2*360)});
  \draw[dashed] (0,1.4) -- (1.55,1.4);
  \node[left] at (0,1.4) {$U_{\max}$};
  \draw (1.55,0.1) -- (1.55,-0.1) node[below] {$\pi/2$};
  \draw[densely dashed] (1.55,0) -- (1.55,1.4);
  \fill (1.05,{1.4*sin(1.05/6.2*360)}) circle (2pt);
  \fill (1.25,{1.4*sin(1.25/6.2*360)}) circle (2pt);
  \node[ForestGreen,right] at (7.0,1.2) {fázová rovnováha:\ $\varphi_0 < \frac{\pi}{2}$};
\end{tikzpicture}
\end{center}

Synchronní částice prochází mezerou ve fázi \(\varphi_0\) na \emph{rostoucí} části sinusovky
(\(\varphi_0<\pi/2\)). Částice, která je pomalejší a přiletí později, pak ,,vidí`` větší
napětí, dostane víc energie a dožene synchronní částici; částice rychlejší přiletí dřív,
dostane méně energie a zpomalí se vůči ostatním. Částice tak oscilují kolem synchronní fáze
(\kw{fázová rovnováha}, autofázování). Kdybychom zvolili \(\varphi_0=\pi/2\) (maximum), získali
bychom sice nejvíc energie, ale svazek by se postupně rozpadl.
\pozor{Pozor:} pro ultrarelativistické částice se rychlost s energií téměř nemění, a v
kruhových urychlovačích (synchrotronech nad tzv. přechodovou energií) je stabilní fáze naopak
na klesající části.

\textbf{Moderní urychlující struktury:}
\begin{itemize}
  \item \kw{Urychlovač s nosnou vlnou}\,\pozn{čtení ,,s nosnou vlnou`` nejisté} --
        \emph{sekundární svazek} budí elektrické pole \(\vec E\parallel\) osa svazku (CLIC).
        V konceptu CLIC (Compact Linear Collider, CERN) intenzivní nízkoenergetický ,,budicí``
        svazek prochází zpomalovacími strukturami, kde vyzáří vysokofrekvenční výkon, a ten se
        převede do urychlujících struktur hlavního svazku (two-beam acceleration).
  \item \kw{Dutinové rezonátory} -- příčná magnetická vlna (\(\vec B\perp\) osa svazku) generuje
        \(\vec E\parallel\) osa svazku. V kovové dutině se vybudí stojaté elektromagnetické vlnění
        (mód \(\mathrm{TM}_{010}\)): magnetické siločáry obíhají kolem osy, elektrické pole míří
        podél osy a urychluje. Supravodivé dutiny dosahují gradientů \(\sim 30\,\mathrm{MV/m}\).
  \item \kw{Laser-wake field}, \kw{plasma-wake field} -- laserový puls nebo svazek částic
        letící plazmatem odtlačí elektrony a za sebou vytvoří ,,brázdu`` (wake) s obrovským
        podélným polem (až \(\sim 100\,\mathrm{GV/m}\), tisíckrát víc než u kovových dutin).
        Perspektivní cesta ke kompaktním urychlovačům.
\end{itemize}

\subsection{Kruhové urychlovače}

Lineární urychlovač využije každou urychlující strukturu jen jednou. Zahneme-li dráhu
magnetickým polem do kruhu, může částice procházet stejnou urychlující mezerou mnohokrát.

\subsubsection{Cyklotron}

\begin{center}
\begin{tikzpicture}[>=Stealth,scale=0.9]
  % duanty
  \draw[thick] (-2.2,0.2) arc[start angle=180,end angle=0,radius=2.2] -- cycle;
  \draw[thick] (-2.2,-0.2) arc[start angle=180,end angle=360,radius=2.2] -- cycle;
  % spirala
  \draw[ForestGreen,thick,domain=0:1800,samples=300,smooth]
     plot ({(0.25+0.105*\x/180)*cos(\x)},{(0.25+0.105*\x/180)*sin(\x)});
  \draw[ForestGreen,thick,->] (1.3,-0.2) .. controls (2.2,0.2) .. (2.9,1.6);
  \fill (0,0) circle (2.5pt);
  \node at (0.5,-1.5) {$\otimes\,\vec B$};
  % RF
  \draw[thick] (-5.0,-0.4) rectangle (-3.8,0.4);
  \draw (-4.4,0) circle (0.25); \node at (-4.4,0) {\small $\sim$};
  \node[left] at (-5.0,0) {RF};
  \draw[thick] (-4.4,0.4) -- (-4.4,1.2) -- (-1.9,1.2);
  \draw[thick] (-4.4,-0.4) -- (-4.4,-1.2) -- (-1.9,-1.2);
\end{tikzpicture}
\end{center}

Cyklotron tvoří dvě duté polokruhové elektrody (,,duanty``, tvaru D) v homogenním
magnetickém poli \(\vec B\) kolmém na rovinu oběhu. Iontový zdroj je uprostřed. V mezeře
mezi duanty je střídavé napětí z RF generátoru; uvnitř duantu je pole nulové a částice
jen opisuje půlkružnici. Při každém průchodu mezerou se urychlí, takže se poloměr dráhy
zvětšuje -- dráha je spirála -- až částice na okraji vyletí ven.

Klíčové je, že \emph{doba oběhu nezávisí na poloměru}. Z rovnosti Lorentzovy a dostředivé
síly \(QevB = mv^2/r\) plyne \(v = QeBr/m\), a tedy oběžná frekvence
\[
  f = \frac{v}{2\pi r} = \frac{QeBr}{2\pi r\, m} = \frac{QeB}{2\pi m} .
\]
Tato \kw{cyklotronová podmínka} umožňuje použít RF generátor o \emph{konstantní} frekvenci
\(f\) -- částice je vždy v mezeře ve správné fázi, ať je na kterémkoli poloměru. Cyklotron
proto může pracovat kontinuálně.

Maximální kinetická energie je dána poloměrem \(R\) magnetu:
\[
  \vzorec{E_k = \tfrac12 mv^2 = \tfrac12 m\,(2\pi R f)^2 = \frac{(QeBR)^2}{2m}}
\]

\begin{priklad}[Odhad]
  Proton (\(Q=1\)) v poli \(B=1{,}5\,\mathrm{T}\) a s poloměrem \(R=0{,}5\,\mathrm{m}\):
  \(f = eB/(2\pi m_p) \approx 23\,\mathrm{MHz}\) a \(E_k = (eBR)^2/(2m_p)\); v praktických
  jednotkách \(pc = 0{,}3\cdot1{,}5\cdot0{,}5\,\mathrm{GeV}=225\,\mathrm{MeV}\), takže
  \(E_k\approx (pc)^2/(2m_pc^2) = 225^2/1876\,\mathrm{MeV}\approx 27\,\mathrm{MeV}\).
\end{priklad}

\strana{P11-7}
\textbf{Relativita:} \(m\to\gamma m\). S rostoucí energií roste relativistická hmotnost a
frekvence oběhu klesá:
\[
  f = \frac{QeB}{2\pi\gamma m} .
\]
Částice se pak začne za RF napětím opožďovat a přestane být urychlována. U protonů to omezuje
klasický cyklotron na \(\sim 20\)--\(30\,\mathrm{MeV}\) (už při \(E_k = 20\,\mathrm{MeV}\) je
\(\gamma\approx1{,}02\), tedy posun o \(2\,\%\)). Řešení jsou dvě:
\begin{enumerate}[label=\alph*)]
  \item \(\vec B\) roste s \(r\) tak, aby \(\gamma(r)/B(r)\) bylo konstantní (\kw{izochronní
        cyklotron}), ale problém s fokusací -- pole rostoucí s poloměrem svazek vertikálně
        rozfokusovává; řeší se to azimutálně proměnným polem (sektorový cyklotron);
  \item frekvenci měnit s energií, \(f\sim\frac1\gamma\) \(\Rightarrow\) \kw{synchrocyklotron}.
        Pak ale lze urychlovat jen jeden shluk částic po druhém -- urychlovač pracuje
        v pulsním režimu (viz Davídkovy slajdy ,,Urychlovače (2)``).
\end{enumerate}
Dalším krokem (pevný poloměr, rostoucí pole i frekvence) je synchrotron, kterým se dosahuje
nejvyšších energií.
